\begingroup
\shorthandoff{"<>} % Previene errores con babel y caracteres especiales en \texttt
\setlength{\arrayrulewidth}{1pt}
\arrayrulecolor{vlepsbased}
\begin{longtable}{|p{15cm}|}
  
  % --- CABECERA INICIAL ---
  \hline
  \rowcolor{maincolor} 
  \multicolumn{1}{|l|}{\textcolor{white}{\textbf{Caso de Uso}}} \\
  \hline
  \endfirsthead
  
  % --- CABECERA EN PÁGINAS SIGUIENTES ---
  \multicolumn{1}{c}{{\tablename\ \thetable{} -- Continuación}} \\
  \hline
  \rowcolor{maincolor} 
  \multicolumn{1}{|l|}{\textcolor{white}{\textbf{Caso de Uso (Continuación)}}} \\
  \hline
  \endhead
  
  % --- PIE DE PÁGINA INTERMEDIO ---
  \hline \multicolumn{1}{|r|}{\textit{Continúa en la siguiente página...}} \\ \hline
  \endfoot
  
  % --- PIE FINAL (Caption abajo) ---
  \hline
  \caption{Caso de Uso CU2 — Compilación y Ejecución de Código en Piston.}
  \label{TB:CU2_PISTON}
  \endlastfoot

  % Fila 1: ID y Actor
  \rowcolor{ulepsbased}
  \begin{tabularx}{\linewidth}{@{}p{7.5cm}p{7.5cm}@{}}
    \textbf{Identificador:} CU2 & \textbf{Actor principal:} Alumno 
  \end{tabularx} \\ \hline

  % Fila 2: Nombre
  \textbf{Nombre:} Compilación y ejecución de código en sandbox \\ \hline

  % Fila 3: Descripción
  \rowcolor{ulepsbased}
  \textbf{Descripción:} \newline 
  El alumno ejecuta el código de su editor en un entorno sandbox aislado (Piston). El sistema envía el código fuente, el lenguaje y la versión al contenedor, y devuelve la salida estándar (\texttt{stdout}), los errores (\texttt{stderr}) y el código de salida. \\ \hline

  % Fila 4: Precondiciones
  \textbf{Precondiciones:} \newline
  -- El alumno ha iniciado sesión en la plataforma. \newline
  -- El contenedor Piston (\texttt{tutor\_compiler}) está activo. \newline
  -- El lenguaje solicitado está instalado en el contenedor Piston. \\ \hline

  % Fila 5: Flujo principal
  \rowcolor{ulepsbased}
  \textbf{Flujo principal:} \newline
  1. El alumno pulsa el botón ``Ejecutar'' en el editor con código en el editor y el lenguaje seleccionado.\newline
  2. El sistema valida los datos de entrada .\newline
  3. El sistema construye el payload y lo envía a la API de Piston.\newline
  4. El sistema envía \texttt{POST} a \texttt{\{PISTON\_URL\}/api/v2/execute}. Timeout de 30 segundos.\newline
  5. Piston ejecuta el código en un sub-contenedor aislado (compila primero si es necesario).\newline
  6. Piston devuelve la respuesta JSON con \texttt{stdout}, \texttt{stderr} y \texttt{exit\_code}.\newline
  7. El sistema procesa la respuesta: concatena errores de compilación y ejecución si existen.\newline
  8. El sistema devuelve al frontend el objeto JSON con la salida procesada y HTTP 200.\newline
  9. El frontend muestra la salida en el panel correspondiente; si hay errores, los resalta. \\ \hline

  % Fila 6: Flujos alternativos
  \textbf{Flujos alternativos:} \newline
  \textbf{FA-01:} Piston no accesible (\texttt{ConnectionError}): Se devuelve HTTP 503 con mensaje de error de servicio.\newline
  \textbf{FA-02:} Timeout de ejecución: Se devuelve HTTP 504 informando que el servicio tardó demasiado.\newline
  \textbf{FA-03:} Señal \texttt{SIGKILL} recibida: El programa excedió el límite de tiempo. Se añade aviso de posible bucle infinito al \texttt{stderr}.\newline
  \textbf{FA-04:} Error de Piston: Se devuelve HTTP 400 con el mensaje específico (ej. lenguaje no soportado).\newline
  \textbf{FA-06:} Error de compilación: Se capturan los errores del compilador y se envían al alumno en el campo \texttt{stderr}. \\ \hline

  % Fila 7: Postcondiciones
  \rowcolor{ulepsbased}
  \textbf{Postcondiciones:} \newline
  1. El alumno recibe la salida de ejecución detallada. \newline
  2. Operación efímera: no se persiste el código ni el resultado en la base de datos. \newline
  3. Los archivos temporales del contenedor se eliminan automáticamente. \\ \hline

\end{longtable}
\endgroup